\documentclass [aspectratio=169]{beamer}
\usetheme{Boadilla}
\usepackage{textpos} % package for the positioning
\usepackage[]{graphicx}
\usepackage{graphicx}
\usepackage{float}
\usepackage{hyperref}
\usepackage{caption}
\usepackage{subcaption}
\usepackage{algorithm,algpseudocode}
\usepackage{amsmath, amsfonts, amsthm, amssymb}
\usepackage{dsfont}
\usepackage[export]{adjustbox}
\usepackage{tikz}
\usetikzlibrary{positioning}
\usetikzlibrary{arrows, shapes, decorations, automata, backgrounds, fit, petri, calc}
\setbeamertemplate{itemize items}[circle]
\setbeamertemplate{enumerate items}[circle]
\usepackage{tikz}
\usepackage{accents}

\newcommand*{\logofont}{\fontfamily{phv}\selectfont}
\definecolor{uoftblue}{RGB}{0,42,92} % official blue color for uoft
\definecolor{deptgreen}{RGB}{114,192,148} 
\definecolor{deptoran}{RGB}{252,103,63} 

\hypersetup{
  colorlinks   = true, %Colours links instead of boxes
  urlcolor     = uoftblue, %Colour for external hyperlinks
  linkcolor    = black, %Colour of internal links
  citecolor   = black %Colour of citations
}

% lin alg
\newcommand{\bu}{{\mathbf{u}}}
\newcommand{\bv}{{\mathbf{v}}}
\newcommand{\bw}{{\mathbf{w}}}
\newcommand{\bx}{\mathbf{x}}
\newcommand{\zerovec}{{\mathbf{0}}}
\newcommand{\inv}{{-1}}


% other useful stuff
\newcommand{\Id}{{\mathds{1}}}
\newcommand{\R}{{\mathbb{R}}}
\newcommand{\C}{{\mathbb{C}}}
\newcommand{\F}{{\mathbb{F}}}
\newcommand{\Z}{{\mathbb{Z}}}
\newcommand{\N}{{\mathbb{N}}}
\newcommand{\Q}{{\mathbb{Q}}}
\newcommand{\cP}{\mathcal{P}}

% set theory
\newcommand{\interior}{\accentset{\circ}}

\newtheorem{proposition}[theorem]{Proposition}


\beamertemplatenavigationsymbolsempty

% block
% example block
% alert block


\title[]{Module 4: Metric Spaces II\\ {\large Operational math bootcamp}\\ \includegraphics[width=8cm]{dept_logo.png}\vspace{-1em}}
\author[]{Ichiro Hashimoto}
\institute[]{University of Toronto}
\date{July 12, 2024}

% set color
\setbeamercolor{title in head/foot}{bg=white}
\setbeamercolor{author in head/foot}{bg=white}
\setbeamercolor{date in head/foot}{fg=uoftblue}
\setbeamercolor{date in head/foot}{bg=white}
\setbeamercolor{title}{fg=uoftblue}
\setbeamerfont{title}{series=\bfseries}
\setbeamercolor{frametitle}{fg=uoftblue}
\setbeamerfont{frametitle}{series=\bfseries}
\setbeamercolor*{item}{fg=uoftblue}
\setbeamercolor{block title}{bg=uoftblue}
\setbeamercolor{block title}{fg=white}
\setbeamercolor{block body}{bg=uoftblue!9!white}
\setbeamercolor{block title example}{bg=deptgreen}
\setbeamercolor{block title example}{fg=white}
\setbeamercolor{block body example}{bg=deptgreen!13!white}
\setbeamercolor{block title alerted}{bg=deptoran}
\setbeamercolor{block title alerted}{fg=white}
\setbeamercolor{block body alerted}{bg=deptoran!10!white}


% set logo at non-title pages
\logo{\includegraphics[height=0.8cm]{dept_logo.png}\vspace*{-.045\paperheight}\hspace*{.78\paperwidth}}

% set margin
\setbeamersize{text margin left=10mm,text margin right=10mm}

\begin{document}
{
\setbeamertemplate{logo}{}
\begin{frame}
    %\vspace{0.5in}
    \titlepage
    %\begin{textblock*}{10cm}(3.5cm,-7.5cm)
      %  \includegraphics[width=8cm]{dept_logo.png}
    %\end{textblock*}
\end{frame}
}

\begin{frame}{Outline}
    \begin{itemize}
      \setlength\itemsep{1em}
      	\item Open and closed sets
    	\item Sequences
	\begin{itemize}
	\item Cauchy sequences
	\item subsequences
	\end{itemize}    
	\end{itemize}
\end{frame}

\begin{frame}
\begin{definition}[Open and closed sets]
Let $(X,d)$ be a metric space. 
\begin{itemize}
    \item A set $U \subseteq X$ is \emph{open} if for every $x \in U$ there exists $\epsilon>0$ such that $B_\epsilon(x) \subseteq U$.
    \item A set $F \subseteq X$ is \emph{closed} if $F^c:= X\setminus F$ is open.
\end{itemize}
\end{definition}

\vspace{0.5em}
\textbf{Note:}


\vspace{1.5em}

\begin{proposition}
Let $(X,d)$ be a metric space. 
 \begin{enumerate}
     \item Let $A_1,A_2\subseteq X$. If $A_1$ and $A_2$ are open, then $A_1 \cap A_2$ is open.
     \item If $A_i \subseteq X$, $i \in I$ are open, then $\cup_{i\in I} A_i$ is open.
 \end{enumerate}
\end{proposition}

\end{frame}


\begin{frame}
\textit{Proof}.
(1) Let $A_1,A_2\subseteq X$. If $A_1$ and $A_2$ are open, then $A_1 \cap A_2$ is open.

\vspace{4cm}

(2) If $A_i \subseteq X$, $i \in I$ are open, then $\cup_{i\in I} A_i$ is open.

\vspace{2cm}

\end{frame}


\begin{frame}
Using DeMorgan, we immediately have the following corollary:
\begin{corollary}
Let $(X,d)$ be a metric space. 
 \begin{enumerate}
     \item Let $A_1,A_2\subseteq X$. If $A_1$ and $A_2$ are closed, then $A_1 \cup A_2$ is closed.
     \item If $A_i \subseteq X$, $i \in I$ are closed, then $\cap_{i\in I} A_i$ is closed.
 \end{enumerate}
\end{corollary}

\end{frame}


\begin{frame}
\begin{definition}[Interior and closure]
Let $A\subseteq X$ where $(X,d)$ is a metric space. 
\begin{itemize}
\setlength\itemsep{1em}
    \item The \emph{closure} of $A$ is $\overline A :=$
    \item The \emph{interior} of $A$ is  $\interior A :=$
    \item The \emph{boundary} of $A$ is $\partial A :=$
\end{itemize}
\vspace{2.5em}
\end{definition}

\vspace{0.3em}

\begin{example}
Let $X = (a, b] \subseteq \R$ with the ordinary (Euclidean) metric. Then

\vspace{1cm}
\end{example}
\end{frame}


\begin{frame}
\begin{proposition}
 Let $A\subseteq X$ where $(X,d)$ is a metric space. Then $\interior A = A \setminus \partial A$.
\end{proposition}
\textit{Proof}.
\vspace{5cm}

\end{frame}

\begin{frame}
\begin{proposition}
 Let $(X,d)$ be a metric space and $A\subseteq X$. $\overline A$ is closed and $\interior A$ is open.
\end{proposition}
\textit{Proof}.
\vspace{9em}

\begin{alertblock}{Remark}

In fact, $\interior A =\bigcup \{U: U \text{ is open and } U \subseteq A \}$ and $\overline{A} = \bigcap \{F \colon F \text{ is closed and } A \subseteq F\}$.
\end{alertblock}

\end{frame}


\begin{frame}{Sequences}
\begin{definition}[Sequence]
Let $(X,d)$ be a metric space. A \emph{sequence} is an ordered list of points $x_n$, $n\in\N$, in $X$, denoted $(x_n)_{n \in \N}$. We say that a sequence $(x_n)_{n \in \N}$ \emph{converges} to a point $x \in X$ if 
\vspace{1.5cm}
%\begin{equation*}
%    \forall \epsilon > 0 \, \exists \, n_\epsilon \in \N \text{ s.t. } d(x_n,x) < \epsilon \text{ for all } n \geq n_\epsilon .
%\end{equation*}
\end{definition}
\end{frame}


\begin{frame}
\textbf{Recall:} $\overline A=$
\vspace{0.25em}

\begin{exampleblock}{Proposition}
Let $(X, d)$ be a metric space, and let $A \subseteq X$. Then $\overline A$ is equal to the set of points in $X$ which are limits of a sequence in $A$.
\end{exampleblock}

\textit{Proof}.
\vspace{4cm}

\end{frame}


\begin{frame}

\vspace{3.25cm}

\begin{alertblock}{Corollary}
\label{cor:closed_converge}
A set $F \subseteq X$, where $(X, d)$ is a metric space, is closed if and only if every sequence in $F$ which converges in $X$ converges to a point in $F$.
\end{alertblock}
\textbf{Remark:}
\end{frame}


\begin{frame}{Cluster points of a set}
\begin{definition}
Let $(X,d)$ be a metric space and $A \subseteq X$. A point $x \in X$ is a \emph{cluster point} of $A$ (also called accumulation point) if for every $\epsilon >0$, $B_\epsilon(x)$ contains infinitely many points in $A$.
\end{definition}
\end{frame}

\begin{frame}

\begin{exampleblock}{Proposition}
 $x \in X$ is a cluster point of $A \subseteq X$ where $(X,d)$ is a metric space if and only if there exists a sequence of points $x_n \in A$, $n \in \N$, such that $x_n \to x$.
\end{exampleblock}
\textit{Proof}.
\vspace{5cm}

\end{frame}


\begin{frame}
Combining the previous result with the limit characterization of closure gives the following: 

\begin{alertblock}{Corollary}
For $A \subseteq X$, $(X,d)$ a metric space, we have $$\overline{A} = A \cup \{x \in X : x \text{ is a cluster point of }A \}.$$
\end{alertblock}

\end{frame}


\begin{frame}{Cauchy sequences}
\begin{definition}[Cauchy sequence]
Let $(X,d)$ be a metric space. A sequence denoted $(x_n)_{n \in \N} \in X$ is called a \emph{Cauchy sequence} if
\vspace{1.5cm}
%\begin{equation*}
%    \forall \epsilon \, \exists \, n_\epsilon \in \N \text{ s.t. } d(x_n,x_m) < \epsilon \text{ for all } n,m \geq n_\epsilon .
%\end{equation*}

\end{definition}

\end{frame}


\begin{frame}
\begin{exampleblock}{Proposition}
\label{prop:converge_means_Cauchy}
Let $(X, d)$ be a metric space, and let $(x_n)_{n\in\N}$ be a convergent sequence in $X$. Then  $(x_n)_{n\in\N}$ is Cauchy.
\end{exampleblock}

\textit{Proof}.
\vspace{4.5cm}
\end{frame}


\begin{frame}
\begin{definition}
A metric space where every Cauchy sequence converges (to a point in the space) is called \emph{complete}.
\end{definition}
\vspace{1em}

\textbf{Example:}

\vspace{2em}

\begin{exampleblock}{Proposition}
Let $(X, d)$ be a metric space, and let $Y\subseteq X$.
\begin{enumerate}
    \item[(i)] If $X$ is complete and if $Y$ is closed in $X$, then $Y$ is complete.
    \item[(ii)] If $Y$ is complete, then it is closed in $X$.
\end{enumerate}
\end{exampleblock}

\end{frame}


\begin{frame}
\textit{Proof}.
\vspace{7cm}
\end{frame}


\begin{frame}{Subsequences}
\begin{definition}
Let $(x_n)_{n \in \N}$ be a sequence in a metric space $(X,d)$. Let $(n_k)_{k \in \N}$ be a sequence of natural numbers with $n_1 < n_2 < \cdots$. The sequence $(x_{n_k})_{k \in \N}$ is called a \emph{subsequence} of $(x_n)_{n \in \N}$. If $(x_{n_k})_{k \in \N}$ converges to $x \in X$, we call $x$ a \emph{subsequential limit}.
\end{definition}

\vspace{0.5cm}

\begin{example}
$\left((-1)^n\right)_{n \in \N}$ 

\vspace{2cm}

%The sequence $\left((-1)^n\right)_{n \in \N}$ diverges but the subsequences $\left((-1)^{2n}\right)_{n \in \N}$ and $\left((-1)^{2n-1}\right)_{n\in \N}$ converge to subsequential limits 1 and $-1$, respectively.
\end{example}

\end{frame}


\begin{frame}
\begin{exampleblock}{Proposition}
A sequence $(x_n)_{n \in \N}$ in a metric space $(X,d)$ converges to $x \in X$ if and only if every subsequence of $(x_n)_{n \in \N}$  also converges to $x$.
\end{exampleblock}
\textit{Proof}.
\vspace{3cm}

\end{frame}


\begin{frame}
\textit{Proof continued}
\vspace{7cm}


\end{frame}

\begin{frame}{References}
 
Charles C. Pugh (2015). \textit{Real Mathematical Analysis}. Undergraduate Texts in Mathematics. \href{https://link-springer-com.myaccess.library.utoronto.ca/book/10.1007/978-3-319-17771-7}{https://link-springer-com.myaccess.library.utoronto.ca/book/10.1007/978-3-319-17771-7}

\vspace{1em}

Runde ,Volker (2005). \textit{A Taste of Topology}. Universitext.  url:  \href{https://link.springer.com/book/10.1007/0-387-28387-0}{https://link.springer.com/book/10.1007/0-387-28387-0} 

\vspace{1em}

Zwiernik, Piotr (2022). \textit{Lecture notes in Mathematics for Economics and Statistics}. url: \href{http://84.89.132.1/~piotr/docs/RealAnalysisNotes.pdf}{http://84.89.132.1/~piotr/docs/RealAnalysisNotes.pdf} \\

\end{frame}




\end{document}